\section{从相位平均理解条纹能否观察}
写下光程差只能确定理想相干分量的明暗位置。实际探测器还会对曝光时间、光谱宽度和光源的空间范围平均；只有这些平均之后的振荡项没有被抵消，条纹才具有可见反差。
沿用前文 $\Gamma_{12}=\langle U_1^*U_2\rangle$ 的定义，它已经包含两路的全部相位。若把可控的传播相位差单独记为 $\Delta\varphi$，就须另外定义去掉这一相位后的复相干度 $\gamma_{12}^{(0)}$，使
\[
 \Gamma_{12}=\sqrt{I_1I_2}\,\gamma_{12}^{(0)}e^{\mathrm i\Delta\varphi},\qquad
 \gamma_{12}=\gamma_{12}^{(0)}e^{\mathrm i\Delta\varphi}.
\]
因此平均强度为
\[
 I=I_1+I_2+2\sqrt{I_1I_2}\,|\gamma_{12}^{(0)}|
 \cos\!\left(\Delta\varphi+\arg\gamma_{12}^{(0)}\right).
\]
两种复相干度的模相同；不能把已包含传播相位的 $\arg\gamma_{12}$ 再与同一 $\Delta\varphi$ 相加。
于是可见度是
\[
 \mathcal V=\frac{I_{\max}-I_{\min}}{I_{\max}+I_{\min}}
 =\frac{2\sqrt{I_1I_2}}{I_1+I_2}|\gamma_{12}|.
\]
即使完全相干，严重不等强的两束光也只有很低的可见度。反过来，平均强度相等也不能保证相干：独立源点之间若没有固定相位，应叠加强度，各自条纹的相位平移会冲淡总图样。

有限谱宽造成的是时间相干性的限制。若两条光路相差 $\Delta\mathcal L$，相应时间延迟为 $\Delta t=\Delta\mathcal L/c$；不同频率分量的相位差随频率改变。当 $\Delta\nu\,|\Delta t|$ 达到一的量级，各频率的条纹不再对齐，反差通常明显下降。因此常用的相干长度量级为 $\ell_c\sim c/\Delta\nu$，但精确数值和衰减曲线取决于谱线形状与采用的宽度定义。

这些限制解释了白光条纹为什么主要出现在近零光程差区域，也解释了为什么窄光谱并不能补救任意宽的空间非相干光源。做题时应先确定几何光程和反射相位，再按题给光源模型进行空间或频率平均；把所有实验困难统称为“相位不同”会遗漏具体可计算的约束。
